% !TeX root = ../本科微分几何笔记.tex
% 由主文档通过 \input 载入；请编译主文档。
\setcounter{chapter}{0} % 保留原书章号。
\chapter{章末习题解答}
\label{app:dg-solutions}

\section{欧氏空间与微分工具}
\begin{solve}[习题~\ref{ex:dg-tools-gram} 的解答]
Gram 行列式为 $\norm v^2\norm w^2-\ip vw^2$，由 Cauchy--Schwarz 不等式非负，亦等于 $\norm{v\times w}^2$。等号当且仅当 $v,w$ 线性相关，其中包括一个向量为零的情形。若两个向量均非零，这等价于两者平行。
\end{solve}

\begin{solve}[习题~\ref{ex:dg-tools-immersion} 的解答]
有 $F_u=(-\sin u,\cos u,0)$、$F_v=(0,0,1)$，两者处处线性无关，因此 $F$ 处处为浸入。但 $F(u+2\pi,v)=F(u,v)$，所以不是单射。限制到 $(u_0-\delta,u_0+\delta)\times(v_0-\varepsilon,v_0+\varepsilon)$，其中 $0<\delta<\pi$，可在相应圆弧上选连续角度分支，从而得到光滑逆与参数片。局部良好不能消除整体周期重复。
\end{solve}

\begin{solve}[习题~\ref{ex:dg-tools-selfadjoint} 的解答]
计算得 $\mathbf g\mathbf A=\begin{pmatrix}0&4\\4&0\end{pmatrix}=\mathbf A^{\mathsf T}\mathbf g$。特征值为 $2,-2$，对应特征向量可取 $(1,2)$、$(1,-2)$。二者关于 $\mathbf g$ 的范数都为 $\sqrt8$，且内积为 $4-4=0$，所以单位特征向量为 $(1,2)/\sqrt8$、$(1,-2)/\sqrt8$。普通矩阵不对称，不妨碍它相对于指定内积自伴。
\end{solve}

\begin{solve}[习题~\ref{ex:dg-tools-cone} 的解答]
定义函数 $F=x^2+y^2-z^2$ 的梯度为 $(2x,2y,-2z)$，仅在原点为零，故除顶点外是正则水平集。若顶点为正则曲面点，通过该点的所有曲面曲线速度都应位于同一个二维切空间内。但 $\gamma_j(t)=tv_j$，其中 $v_1=(1,0,1),v_2=(0,1,1),v_3=(-1,0,1)$，都完全位于圆锥上，而三个速度线性无关，矛盾。这里使用了切空间的几何性质，而非仅以梯度判据失效作结论。
\end{solve}

\section{空间曲线}
\begin{solve}[习题~\ref{ex:dg-curves-parabola} 的解答]
有 $\gamma'=(1,2t,0)$、$\gamma''=(0,2,0)$，所以 $\kappa(t)=2/(1+4t^2)^{3/2}$。原点处单位切向为 $(1,0,0)$，主法向为 $(0,1,0)$，曲率为 $2$，故密切圆位于 $z=0$ 平面，中心 $(0,1/2,0)$，半径 $1/2$，方程为 $x^2+(y-1/2)^2=1/4$。其下半圆在原点附近有展开 $y=x^2+O(x^4)$，与抛物线的二阶项一致。
\end{solve}

\begin{solve}[习题~\ref{ex:dg-curves-reverse} 的解答]
先取 $\widetilde\gamma(t)=\gamma(-t)$，三个导数依次增加系数 $-1,1,-1$，因此叉积范数不变，三重积系数为 $(-1)(1)(-1)=1$，曲率与挠率均不变。一般反向参数变换还会增加沿低阶导数的项，但它们在叉积和三重积中消失，最终分子分母具有相同的参数缩放因子。对空间正交变换 $Q$，叉积范数不变，而三重积乘以 $\det Q$；故 $\det Q=-1$ 时挠率变号，曲率不变。平移不影响任何导数。
\end{solve}

\begin{solve}[习题~\ref{ex:dg-curves-lancret} 的解答]
若 $\tau/\kappa=m$ 为常数，令 $a=(m\mathbf t+\mathbf b)/\sqrt{1+m^2}$。由 Frenet 方程，$a'=(m\kappa-\tau)\mathbf n/\sqrt{1+m^2}=0$，且 $\ip{\mathbf t}a=m/\sqrt{1+m^2}$ 为常数。反之，设固定单位向量 $a$ 满足 $\ip{\mathbf t}a=c$，求导得 $\ip{\mathbf n}a=0$，故 $a=c\mathbf t+d\mathbf b$。由单位长度，$d^2=1-c^2$；若 $d=0$，则 $\mathbf t$ 固定，与 $\kappa>0$ 矛盾。因此 $d$ 连续且非零，在连通区间上为常数。求导得 $(c\kappa-d\tau)\mathbf n=0$，故 $\tau/\kappa=c/d$。
\end{solve}

\begin{solve}[习题~\ref{ex:dg-curves-circle} 的解答]
由 $\tau=0$ 得曲线在一个平面中，且 $c_0=\gamma+\mathbf n/c$ 的导数为 $\mathbf t-\mathbf t=0$，所以 $c_0$ 为常点。又 $\norm{\gamma-c_0}=1/c$，曲线位于该平面内以 $c_0$ 为中心的圆上。反之，半径 $1/c$ 的圆按弧长参数化后，直接求导得到曲率 $c$ 与挠率 $0$。结论描述曲线的像，未额外排除多次绕行。
\end{solve}

\section{平面曲线}
\begin{solve}[习题~\ref{ex:dg-plane-ellipse} 的解答]
分子 $x'y''-y'x''=ab$，分母为 $(a^2\sin^2t+b^2\cos^2t)^{3/2}$，所以
\[
 k(t)=\frac{ab}{(a^2\sin^2t+b^2\cos^2t)^{3/2}}>0.
\]
当 $a>b$ 时，分母在长轴端点 $t=0,\pi$ 最小，故最大曲率为 $a/b^2$；在短轴端点 $t=\pi/2,3\pi/2$ 最小曲率为 $b/a^2$。长轴端点弯得更急，与曲率半径较小相一致。
\end{solve}

\begin{solve}[习题~\ref{ex:dg-plane-closed} 的解答]
设 $\Theta(s)=\int_0^s k(\sigma)\dd\sigma$。充分必要条件为
\[
 \Theta(L)\in2\pi\mathbb Z,\qquad
 \int_0^L\cos\Theta(s)\dd s=0,\qquad
 \int_0^L\sin\Theta(s)\dd s=0.
\]
第一条件保证切向量首尾一致，后两条件保证位置闭合；$k$ 的光滑周期性进一步保证所有高阶导数匹配。改变初始角度只把位移积分整体旋转，所以其是否为零不变。这些条件保证正则闭合，但不保证简单闭合。
\end{solve}

\begin{solve}[习题~\ref{ex:dg-plane-offset} 的解答]
由 $n'=-k\mathbf t$ 得 $\gamma_a'=(1-ak)\mathbf t$。在 $1-ak>0$ 的区间内，$\dd s_a=(1-ak)\dd s$，单位切向仍为 $\mathbf t$，于是 $k_a=\ip{\dd\mathbf t/\dd s_a}n=k/(1-ak)$。若 $1-ak=0$，偏移参数曲线速度为零，正则性失效，不能继续应用曲率公式。常曲率圆在向内偏移一个曲率半径时整圈缩为圆心，是这种退化的直接例子。
\end{solve}

\begin{solve}[习题~\ref{ex:dg-plane-area} 的解答]
Green 公式给出 $A=\frac12\int_0^{2\pi}(xy'-yx')\dd t=\pi ab$。由等周不等式，$L\ge2\pi\sqrt{ab}$，等号当且仅当椭圆为圆，即 $a=b$。这给出可直接计算的下界，并没有给出一般椭圆周长的初等闭式。
\end{solve}

\section{正则曲面}
\begin{solve}[习题~\ref{ex:dg-surfaces-tangent} 的解答]
梯度在椭球上非零，切向量满足 $x_0X/a^2+y_0Y/b^2+z_0Z/c^2=0$。仿射切平面为
\[
 \frac{x_0(x-x_0)}{a^2}+\frac{y_0(y-y_0)}{b^2}
 +\frac{z_0(z-z_0)}{c^2}=0,
\]
利用点在椭球上，可化为 $x_0x/a^2+y_0y/b^2+z_0z/c^2=1$。
\end{solve}

\begin{solve}[习题~\ref{ex:dg-surfaces-chart} 的解答]
有 $\mathbf x_\theta\times\mathbf x_\varphi=R^2\sin\theta(\sin\theta\cos\varphi,\sin\theta\sin\varphi,\cos\theta)$，故 $\theta=0,\pi$ 处退化。北极附近可取 $\mathbf y(u,v)=(u,v,\sqrt{R^2-u^2-v^2})$，定义域为 $u^2+v^2<R^2$；前两个坐标投影为逆映射，叉积第三分量为 $1$，因此这是覆盖北极的正则参数片。经度还需限制在长度小于 $2\pi$ 的开区间内以消除重复。
\end{solve}

\begin{solve}[习题~\ref{ex:dg-surfaces-orientation} 的解答]
两个单位法向量都位于同一条法线，所以 $\widetilde\nu(p)=\sigma(p)\nu(p)$，其中 $\sigma(p)=\ip{\widetilde\nu(p)}{\nu(p)}\in\{1,-1\}$ 连续。连通性使 $\sigma$ 恒定。若曲面是两个互不相交的球面的并，可在第一个分支保持法向、在第二个分支反向，得到既不处处相同也不处处相反的两个法向场。
\end{solve}

\begin{solve}[习题~\ref{ex:dg-surfaces-differential} 的解答]
微分为 $\dd F_{(u,v)}(a,b)=(a,b,2ua+2vb)$，它是从平面到图像切空间的线性同构。映射 $F$ 双射，其逆为 $(x,y,z)\mapsto(x,y)$ 在抛物面上的限制，两者都光滑，因此为微分同胚。但它通常不保持长度：像向量的范数平方为 $a^2+b^2+(2ua+2vb)^2$，一般不等于 $a^2+b^2$。
\end{solve}

\section{第一基本形式}
\begin{solve}[习题~\ref{ex:dg-metric-graph} 的解答]
有 $E=1+f_u^2$、$F=f_uf_v$、$G=1+f_v^2$，故 $EG-F^2=1+f_u^2+f_v^2$，面积元为 $\sqrt{1+f_u^2+f_v^2}\dd u\dd v$。当 $f=uv$ 时，$\mathrm I=(1+v^2)\dd u^2+2uv\dd u\dd v+(1+u^2)\dd v^2$，面积元为 $\sqrt{1+u^2+v^2}\dd u\dd v$。应保留混合项，否则会错误计算非坐标方向的长度。
\end{solve}

\begin{solve}[习题~\ref{ex:dg-metric-sphere} 的解答]
度量为 $R^2\dd\theta^2+R^2\sin^2\theta\dd\varphi^2$。固定 $\theta=\theta_0$ 的纬线长度为 $2\pi R\sin\theta_0$，面积元为 $R^2\sin\theta\dd\theta\dd\varphi$。在 $0<\theta<\pi$、$0<\varphi<2\pi$ 积分得面积 $4\pi R^2$；缺少的经线与两极为面积零集合，可通过取极限或增加坐标片处理。
\end{solve}

\begin{solve}[习题~\ref{ex:dg-metric-cylinder} 的解答]
任何连接两点的圆柱曲线都可沿区间连续选择角度的实值提升；其终点角度差为 $\Delta\theta+2\pi m$，其中 $m\in\mathbb Z$。在展开平面中，这条曲线的长度保持不变，至少为两提升端点的欧氏距离 $\sqrt{R^2(\Delta\theta+2\pi m)^2+\Delta z^2}$。反过来，每一对提升端点之间的直线都可卷回圆柱，达到该长度。由于表达式在 $|m|\to\infty$ 时趋于无穷，下确界由某个整数达到，遂得公式；角度差为 $\pi$ 模 $2\pi$ 时可能出现两条不同最短路径。
\end{solve}

\begin{solve}[习题~\ref{ex:dg-metric-conformal} 的解答]
取原空间标准正交基 $e_1,e_2$。保角使 $Ae_1\perp Ae_2$，记二者长度为 $a,b>0$。由于 $e_1+e_2$ 与 $e_1-e_2$ 正交，像也正交，故 $a^2-b^2=0$。于是任意向量 $X=x_1e_1+x_2e_2$、$Y=y_1e_1+y_2e_2$ 满足 $\ip{AX}{AY}=a^2(x_1y_1+x_2y_2)=a^2\ip XY$。可逆性排除了缩放因子为零。
\end{solve}

\section{第二基本形式与曲率}
\begin{solve}[习题~\ref{ex:dg-curvature-quadric} 的解答]
原点一阶导数为零，因此第一基本形式是单位矩阵，形算子为 Hessian 矩阵 $\operatorname{diag}(2,-3)$。主曲率为 $2,-3$，$K=-6$、$H=-1/2$。渐近方向满足 $2a^2-3b^2=0$，即 $b=\pm\sqrt{2/3}\,a$。它是双曲点，主方向为两个坐标轴方向。
\end{solve}

\begin{solve}[习题~\ref{ex:dg-curvature-normal} 的解答]
反转法向后 $\widetilde S=-S$，因此每个特征向量仍是特征向量，其特征值变号；$\widetilde H=-H$，$\widetilde K=\det(-S)=\det S=K$，这里使用了切空间维数为 $2$。渐近方向的条件 $\mathrm{II}(X,X)=0$ 在整体乘以 $-1$ 后不变。若按大小重新编号主曲率，编号可能互换，但对应方向集合不变。
\end{solve}

\begin{solve}[习题~\ref{ex:dg-curvature-umbilic} 的解答]
若 $\lambda=0$，则 $\dd\nu=0$，连通性使法向恒定。函数 $\ip{\mathbf x}\nu$ 的微分为零，故曲面包含于一个固定平面。若 $\lambda\ne0$，则 $\dd(\mathbf x+\nu/\lambda)=0$，所以 $c=\mathbf x+\nu/\lambda$ 是常点，且 $\norm{\mathbf x-c}=1/|\lambda|$。曲面因而包含于相应球面。题目只给出包含关系；不加完备或紧致等条件时，它可以仅是球面的一块开片。
\end{solve}

\begin{solve}[习题~\ref{ex:dg-curvature-matrices} 的解答]
由于混合项是 $2f\dd u\dd v$，有 $f=2$，所以 $\mathbf b=\begin{pmatrix}0&2\\2&0\end{pmatrix}$，$[S]=\begin{pmatrix}0&1/2\\2&0\end{pmatrix}$。主曲率为 $1,-1$，对应坐标向量 $(1,2)$、$(1,-2)$，归一化后分别除以 $\sqrt8$。若把 $\mathbf b$ 的普通特征值 $2,-2$ 当作主曲率，就忽略了坐标基的长度差异。正确问题是广义特征值方程 $\mathbf b a=k\mathbf g a$。
\end{solve}

\section{典型曲面}
\begin{solve}[习题~\ref{ex:dg-examples-cone} 的解答]
直接求内积得 $\mathrm I=(1+c^2)\dd r^2+r^2\dd\theta^2$。令 $\rho=\sqrt{1+c^2}\,r$、$\phi=\theta/\sqrt{1+c^2}$，则 $\mathrm I=\dd\rho^2+\rho^2\dd\phi^2$，正是平面极坐标度量。所以 $(r,\theta)\mapsto(\rho\cos\phi,\rho\sin\phi)$ 为局部等距映射。绕锥一周时展开角为 $2\pi/\sqrt{1+c^2}$；完整展开需切开一条母线，顶点仍不在正则讨论范围内。
\end{solve}

\begin{solve}[习题~\ref{ex:dg-examples-torus} 的解答]
面积元为 $a(R+a\cos u)\dd u\dd v$，所以 $K\dd A=\cos u\dd u\dd v$。积分得
\[
 \int_\Sigma K\dd A=\int_0^{2\pi}\int_0^{2\pi}\cos u\dd u\dd v=0.
\]
正曲率区域的积分为 $4\pi$，负曲率区域的积分为 $-4\pi$，二者抵消。这是带符号的曲率积分为零，不是曲率处处为零，也不是绝对曲率积分为零。
\end{solve}

\begin{solve}[习题~\ref{ex:dg-examples-helicoid} 的解答]
有 $\mathbf y_u=(\cos v,\sin v,0)$、$\mathbf y_v=(-u\sin v,u\cos v,a)$，故 $E=1,F=0,G=u^2+a^2$。叉积为 $(a\sin v,-a\cos v,u)$，除以 $\sqrt{u^2+a^2}$ 得单位法向。二阶导数为 $\mathbf y_{uu}=0$、$\mathbf y_{uv}=(-\sin v,\cos v,0)$、$\mathbf y_{vv}=(-u\cos v,-u\sin v,0)$，得到 $e=g=0$、$f=-a/\sqrt{u^2+a^2}$。式~\eqref{eq:dg-KH} 遂给出 $H=0$、$K=-a^2/(u^2+a^2)^2$。
\end{solve}

\begin{solve}[习题~\ref{ex:dg-examples-sphere-variation} 的解答]
直接求导为 $\left.\dd[4\pi(R+\varepsilon)^2]/\dd\varepsilon\right|_0=8\pi R$。外法向下 $H=-1/R$，取 $f=1$，无边界时不需要边界支撑限制，公式给出 $-2\int_\Sigma fH\dd A=(2/R)(4\pi R^2)=8\pi R$。这个检验同时检查了形算子的负号和第一变分的负号。
\end{solve}

\section{协变导数与基本方程}
\begin{solve}[习题~\ref{ex:dg-connection-warped} 的解答]
度量矩阵为 $\operatorname{diag}(1,r^2)$，仅 $\partial_sg_{\theta\theta}=2rr'$ 非零，故
\[
 \Gamma^s_{\theta\theta}=-rr',\qquad
 \Gamma^\theta_{s\theta}=\Gamma^\theta_{\theta s}=r'/r.
\]
曲率分子为 $\partial_s(-rr')-(-rr')(r'/r)=-(r'^2+rr'')+r'^2=-rr''$，分母为 $r^2$，所以 $K=-r''/r$。这是只依赖度量的计算，适用于任何具有此度量的曲面，不要求先给出旋转嵌入。
\end{solve}

\begin{solve}[习题~\ref{ex:dg-connection-conformal} 的解答]
取正交标架 $e_1=e^{-\varphi}\partial_u,e_2=e^{-\varphi}\partial_v$。由正交坐标公式，$\omega=-\varphi_v\dd u+\varphi_u\dd v$，因此 $\dd\omega=(\varphi_{uu}+\varphi_{vv})\dd u\wedge\dd v$。又 $\dd A=e^{2\varphi}\dd u\wedge\dd v$，式~\eqref{eq:dg-connection-curvature} 给出 $K=-e^{-2\varphi}(\varphi_{uu}+\varphi_{vv})$。立体投影度量中 $\varphi=\log2-\log(1+u^2+v^2)$，其欧氏 Laplace 为 $-4/(1+u^2+v^2)^2$，故 $K=1$。
\end{solve}

\begin{solve}[习题~\ref{ex:dg-connection-incompatible} 的解答]
由第一基本形式的常系数可知全部 Christoffel 符号为零，内蕴曲率为零。但给定第二基本形式使 $[S]=I_2$，外蕴公式给出 $K=1$，违背 Gauss 方程。因此不存在这样的曲面片。正定第一基本形式和对称第二基本形式各自合法，不意味着它们能同时来自同一个曲面。
\end{solve}

\begin{solve}[习题~\ref{ex:dg-connection-parallel} 的解答]
平行向量场满足 $\nabla_{\dot\gamma}V=0$，坐标式为
\[
 \dot V^k+\sum_{i,j}\Gamma^k_{ij}(u(t))\dot u^iV^j=0.
\]
这是关于 $V^k$ 的线性方程，系数沿给定曲线光滑，因此给定初值后唯一存在，可跨坐标片依唯一性拼接；在曲线定义域的每个紧子区间上用有限坐标片完成延伸。两个平行场 $V,W$ 满足 $\dd\ip VW/\dd t=\ip{\nabla_{\dot\gamma}V}W+\ip V{\nabla_{\dot\gamma}W}=0$，所以平行移动保持内积、长度和夹角。
\end{solve}

\section{测地线}
\begin{solve}[习题~\ref{ex:dg-geodesics-latitude} 的解答]
纬线空间中是半径 $R\sin\theta$ 的圆，故 $\kappa=1/(R\sin\theta)$。按经度增加方向，$t=(-\sin\varphi,\cos\varphi,0)$，$q=\nu\times t=(-\cos\theta\cos\varphi,-\cos\theta\sin\varphi,\sin\theta)$，且 $t'=-(\cos\varphi,\sin\varphi,0)/(R\sin\theta)$，这里撇表示弧长导数。于是 $k_n=-1/R$、$k_g=\cot\theta/R$，满足 $\kappa^2=k_n^2+k_g^2$。仅在 $\theta=\pi/2$ 时测地曲率为零，即只有赤道是这族纬线中的测地线。
\end{solve}

\begin{solve}[习题~\ref{ex:dg-geodesics-equator} 的解答]
纬线 $s=s_0$ 上 $\dot s=\ddot s=0$，非恒定性保证 $\dot\theta\ne0$。第一条测地线方程成为 $r(s_0)r'(s_0)\dot\theta^2=0$；由于 $r>0$，等价于 $r'(s_0)=0$。反之，若该条件成立，令 $\theta$ 为时间的仿射函数，两方程均满足。标准环面母线以 $s=au$ 为弧长，$r=R+a\cos u$，所以 $r'= -\sin u$，得到外赤道 $u=0$ 与内赤道 $u=\pi$。
\end{solve}

\begin{solve}[习题~\ref{ex:dg-geodesics-reparameter} 的解答]
链式法则和联络乘积法则给出
\[
 \nabla_{\frac{\dd}{\dd t}(\gamma\circ\phi)}
       \frac{\dd}{\dd t}(\gamma\circ\phi)
 =(\phi')^2(\nabla_{\dot\gamma}\dot\gamma)\circ\phi
 +\phi''\dot\gamma\circ\phi
 =\phi''\dot\gamma\circ\phi.
\]
非恒定仿射参数测地线的速度为非零常数，因此该式为零当且仅当 $\phi''=0$，即 $\phi(t)=at+b$，其中微分同胚条件要求 $a\ne0$。
\end{solve}

\begin{solve}[习题~\ref{ex:dg-geodesics-hyperbolic} 的解答]
取 $\varphi=-\log y$，有 $\varphi_{xx}+\varphi_{yy}=1/y^2$，故 $K=-1$。非零 Christoffel 符号为 $\Gamma^x_{xy}=\Gamma^x_{yx}=-1/y$、$\Gamma^y_{xx}=1/y$、$\Gamma^y_{yy}=-1/y$，测地线方程是
\[
 x''-2x'y'/y=0,\qquad y''+(x'^2-y'^2)/y=0.
\]
当 $x'=0$ 时，$y(t)=Ae^{bt}$，$A>0$，给出竖直测地线，$b=0$ 为常值曲线。非竖直时可按单位速参数化，第一式给出 $x'/y^2=c\ne0$，速度条件为 $(x'^2+y'^2)/y^2=1$。第二式进一步说明 $\dd(y'/y)/\dd t=-c^2y^2$，故 $x+y'/(cy)=a$ 为常数。由速度条件消去 $y'$，得到 $(x-a)^2+y^2=1/c^2$。因此非竖直测地线的像为圆心位于边界直线上的上半圆弧；这里“所有”是指最大延伸后整条上半圆，也包括其子弧。
\end{solve}

\section{Gauss--Bonnet 定理}
\begin{solve}[习题~\ref{ex:dg-gb-triangle} 的解答]
角盈为 $3(2\pi/3)-\pi=\pi$。由恒曲率 $K=1/R^2$，面积为 $\pi R^2$。这个数是整个球面积的四分之一；题目指定短大圆弧及其三角形区域，避免把互补区域误当成同一个三角形。
\end{solve}

\begin{solve}[习题~\ref{ex:dg-gb-annulus} 的解答]
取向上法向，外边界按逆时针方向，$k_g=1/b$、长度 $2\pi b$，积分为 $2\pi$；内边界必须顺时针，使圆环区域仍在左侧，所以 $k_g=-1/a$、长度 $2\pi a$，积分为 $-2\pi$。内部 $K=0$，无角点，左端总和为零。圆环可收缩到圆周，欧拉示性数为零，故公式成立；若把内边界也取逆时针，就破坏了边界正向约定。
\end{solve}

\begin{solve}[习题~\ref{ex:dg-gb-positive} 的解答]
由紧致性和 $K>0$，曲率积分严格为正。Gauss--Bonnet 给出 $2\pi(2-2g)>0$，且 $g$ 是非负整数，因此 $g=0$。该结论仅描述拓扑类型。非球形椭球有正高斯曲率且拓扑上是球面，却不具有常曲率，也不是某个标准圆球的刚体运动像。
\end{solve}

\begin{solve}[习题~\ref{ex:dg-gb-flat-embedded} 的解答]
固定 $a\in\R^3$，令 $F(p)=\tfrac12\norm{p-a}^2$，它在紧致曲面上达到最大值，取最大点 $p$，并记 $R=\norm{p-a}>0$。因为没有边界，对任意切向量 $v$ 有 $\ip{p-a}v=0$，所以可选局部法向 $\nu=(p-a)/R$。对任意单位切向量 $v$，选择曲面曲线 $\gamma(0)=p,\gamma'(0)=v$，最大值条件给出
\[
 0\ge(F\circ\gamma)''(0)=1+\ip{p-a}{\gamma''(0)}
 =1+R\mathrm{II}(v,v).
\]
因此所有单位方向法曲率都不大于 $-1/R$，两个主曲率均严格为负，得到 $K(p)\ge1/R^2>0$，矛盾。证明只用到最大点附近的法向，不要求曲面预先给定整体定向。这个结论也说明抽象平坦环面不能以该度量光滑等距嵌入 $\R^3$。
\end{solve}

\section{流形与黎曼度量}
\begin{solve}[习题~\ref{ex:dg-manifolds-gradient} 的解答]
极坐标的逆度量为 $\operatorname{diag}(1,r^{-2})$，所以 $\grad f=f_r\partial_r+(f_\theta/r^2)\partial_\theta$。由 $\sqrt{|g|}=r$，有 $\Delta f=f_{rr}+f_r/r+f_{\theta\theta}/r^2$。对 $f=r^2$，得到 $\grad f=2r\partial_r$、$\Delta f=4$，与直角坐标下 $f=x^2+y^2$ 的计算一致。$r=0$ 的坐标表达失效，但对应函数和向量场可在直角坐标中光滑延伸。
\end{solve}

\begin{solve}[习题~\ref{ex:dg-manifolds-sphere-laplace} 的解答]
球面度量给出
\[
 \Delta f=\frac1{R^2\sin\theta}\partial_\theta(\sin\theta\,\partial_\theta f)
 +\frac1{R^2\sin^2\theta}\partial_\varphi^2f.
\]
取 $f=\cos\theta$ 得 $\Delta f=-2\cos\theta/R^2$。虽然经纬表达在极点含有形式上的零分母，$f=z/R$ 是整个球面上的光滑函数，$\Delta f=-2f/R^2$ 也光滑；在别的坐标中计算或取此光滑延伸即可理解极点处的值。
\end{solve}

\begin{solve}[习题~\ref{ex:dg-manifolds-flat-torus} 的解答]
整数平移是欧氏等距，因此同一商点的不同提升给出同样的切向量内积，度量定义无歧义。商映射在足够小的邻域内为局部等距，测地线方程可提升为 $x''=y''=0$，故所有测地线来自直线 $p+t(a,b)$ 的投影。闭合的充要条件是存在 $T>0$，使 $T(a,b)\in\mathbb Z^2$。若 $ab\ne0$，这等价于 $a/b\in\mathbb Q$；若一个分量为零，另一分量非零，也总能选取周期。速度不能同时为零，已由题目排除。
\end{solve}

\begin{solve}[习题~\ref{ex:dg-manifolds-scale} 的解答]
新长度为原长度的 $c$ 倍，新面积为原面积的 $c^2$ 倍。因为缩放因子恒定，Christoffel 符号中逆度量的 $c^{-2}$ 与度量导数的 $c^2$ 抵消，联络及 $(1,3)$ 型曲率算子保持不变；用新标准正交基 $e_i/c$ 计算高斯曲率，得到 $\widetilde K=K/c^2$。因此 $\widetilde K\dd\widetilde A=K\dd A$，总曲率不变，与 Gauss--Bonnet 的拓扑性质一致。若缩放因子随位置变化，曲率还会含有其导数，不能使用此结论。
\end{solve}

